\subsection{诱导滤子}

命题 8. 设 \(\mathfrak{F}\) 是集合 X 上的滤子，A 是 X 的子集。则迹 \(\mathfrak{F}_{\mathrm{A}}\)（\(\mathfrak{F}\) 在 A 上的）是滤子，当且仅当 \(\mathfrak{F}\) 的每个集合都与 A 相交。

由 \((\mathbf{M} \cap \mathbf{N}) \cap \mathbf{A} = (\mathbf{M} \cap \mathbf{A}) \cap (\mathbf{N} \cap \mathbf{A})\) 可见 \(\mathfrak{F}_{\mathbf{A}}\) 满足 \((\mathbf{F}_{\mathbf{II}})\)；再者，若 \(\mathbf{M} \cap \mathbf{A} \subset \mathbf{P} \subset \mathbf{A}\)，则 \(\mathbf{P} = (\mathbf{M} \cup \mathbf{P}) \cap \mathbf{A}\)，从而 \(\mathfrak{F}_{\mathbf{A}}\) 满足 \((\mathbf{F}_{\mathbf{I}})\)。于是 \(\mathfrak{F}_{\mathbf{A}}\) 是滤子，当且仅当它满足 \((\mathbf{F}_{\mathbf{III}})\)，亦即当且仅当 \(\mathfrak{F}\) 的每个集合都与 \(\mathbf{A}\) 相交。

特别地，若 \(\mathbf{A} \in \mathfrak{F}\)，则 \(\mathfrak{F}_{\mathbf{A}}\) 是 \(\mathbf{A}\) 上的滤子（由 \((\mathbf{F}_{\mathbf{II}})\) 与 \((\mathbf{F}_{\mathbf{III}})\)）。

定义 5. 设 A 是集合 X 的子集，F 是 X 上的滤子。若 F 在 A 上的迹是 A 上的滤子，则称此滤子为 F 在 A 上诱导的滤子。

若滤子 \(\mathfrak{F}\)（\(X\) 上的）在 \(A \subset X\) 上诱导出一个滤子，则在 \(A\) 上，\(\mathfrak{F}\) 的一个基的迹是 \(\mathfrak{F}_A\) 的基，由第 3 目命题 3。

例. 设 X 是拓扑空间，A 是 X 的子集，x 是 X 的一点。为使 x 的邻域滤子 B 在 A 上的迹是 A 上的滤子，必须且只需 x 的每个邻域都与 A 相交，亦即 x 属于 A 的闭包（§ 1，第 6 目，定义 10）。

这个诱导滤子的例子之所以重要，有两个原因：第一，它在极限理论中起重要作用（§ 7，第 5 目）；第二，每个滤子都可以用这种方式定义。事实上，设 F 是集合 X 上的滤子，并设 X' 是在 X 上添加一个新元素 ω 所得的集合，其中 X 与 \{ω\} 在 X' 中的余集视为同一（《集合论》，R，§ 4，第 5 目）；设 F' 是 X' 上由形如 M ∪ \{ω\} 的集合组成的滤子，其中 M 取遍 F。对 X' 的每个点 x ≠ ω，设 B(x) 是 X' 中包含 x 的所有子集所成之集，并设 B(ω) 为 F'；于是 x ∈ X' 的诸 B(x) 显然满足公理 (V₁)、(V₂)、(V₃) 与 (V₄)，从而在 X' 上定义了一个拓扑，使它们成为各点的邻域滤子。最后，在此拓扑下 ω 属于 X 的闭包，且 F 是由 F' = B(ω) 在 X 上诱导的。这样在 X' 上定义的拓扑（相应地，赋以该拓扑的集合 X'）称为伴随于 F 的拓扑（相应地，拓扑空间）。

命题 9. 超滤子 \(\mathfrak{U}\)（集合 \(X\) 上的）在子集 \(A\)（\(X\) 的）上诱导出滤子，当且仅当 \(A \in \mathfrak{U}\)；且当此条件满足时，\(\mathfrak{U}_A\) 是 \(A\) 上的超滤子。

这是第 4 目命题 5 与命题 6 的直接推论。

